% ============================================================
%  Trening Mózgu -- book preamble
%
%  THE UNIT IS A SET, NOT AN EXERCISE.
%
%  A set is a page of ~30 exercises listed one under another, done against a
%  stopwatch and scored as a whole. Answers live in one appendix at the BACK of
%  the book, never beside the exercise.
%
%  This is a deliberate replacement of an earlier design that gave each
%  exercise its own frame, with the answer opening the next one. That design is
%  right for programmed instruction, where the point is to LEARN a step. It is
%  wrong here: this book drills, and what it measures is how many you get
%  through and how fast. One exercise per frame put two on a page, made the
%  page ceremony rather than work, and put the answer where it interrupts the
%  run. Nothing about it is kept.
%
%  Engine: pdfLaTeX. Do NOT load newtx -- the companion repository probes for
%  it and degrades, which made its TS1 font map a trap on a half-installed
%  machine. This preamble simply never asks.
% ============================================================

\usepackage[T1]{fontenc}
\usepackage[utf8]{inputenc}
\usepackage{lmodern}
\usepackage[polish]{babel}

\usepackage{amsmath}
\usepackage{amssymb}
\usepackage{etoolbox}
\usepackage{xcolor}
\usepackage{multicol}
\usepackage{tikz}
\usepackage{fontawesome5}
\usepackage{fancyhdr}
\usepackage[hidelinks]{hyperref}

% ------------------------------------------------------------
%  Paper
% ------------------------------------------------------------
% A4 at 11pt. Margins are TIGHTER than a prose book's on purpose: the page is a
% worksheet, and every millimetre of margin is an exercise that did not fit.
% There is no outer-margin furniture any more, so nothing needs room out there.
\usepackage[
  a4paper,
  inner=1.8cm, outer=1.5cm, top=1.8cm, bottom=1.6cm,
  head=24pt, headsep=0.5cm, footskip=0.9cm
]{geometry}

% ------------------------------------------------------------
%  Colours -- carried over from the archived deck
% ------------------------------------------------------------
\definecolor{ArithColor}{HTML}{6BAED6}
\definecolor{OpsColor}{HTML}{FD8D3C}
\definecolor{LogicColor}{HTML}{9C27B0}
\definecolor{SeqColor}{HTML}{388E3C}
\definecolor{MixColor}{HTML}{D81B60}
\definecolor{TrickColor}{HTML}{00838F}

\definecolor{btrule}{HTML}{C8D0D8}
\definecolor{btink}{HTML}{1F4E79}
\definecolor{btgrey}{HTML}{5A5A5A}

% ------------------------------------------------------------
%  The answer store
% ------------------------------------------------------------
% Answers are written at the point of use and printed at the BACK. Carrying
% them there through a global macro store rather than an .aux file is
% deliberate: material written to an aux-style file is expanded at shipout,
% where a fragile command breaks with an error naming neither the answer nor
% the set it came from. \include is sequential within a run, so a macro defined
% globally while typesetting set 3 is still defined when the appendix is set.
%
% \csxappto expands \thezad at store time -- the number has to be frozen when
% the exercise is written, not when the appendix is -- while \unexpanded keeps
% the answer's own tokens literal, so $ and \frac survive to the back page.
\newcounter{zestaw}
\newcounter{zad}[zestaw]

% \btsep, not \quad. An answer list is a long run of short unbreakable items
% with no hyphenation anywhere in it, so a rigid separator leaves TeX no legal
% breakpoint it can afford and every line comes out overfull -- measured at
% 12.0pt, seventeen times, on the first build of this appendix. The penalty
% offers the breakpoint and the stretch lets the line end early.
\newcommand{\btsep}{\penalty0\hskip 0.9em plus 0.6em minus 0.25em\relax}
\newcommand{\btstoreans}[1]{%
  \csxappto{btans@\thezestaw}{\thezad)~\unexpanded{#1}\noexpand\btsep}%
}

\newlength{\btzadskip}
\setlength{\btzadskip}{3.2pt}

% ------------------------------------------------------------
%  \z -- one exercise on one line
% ------------------------------------------------------------
% A number, the problem, and a rule to write on. That is the whole element.
% There is no difficulty star and no per-exercise timer: both are properties of
% the SET, printed once at its head, and repeating them thirty times a page
% would be noise that costs exercises.
%
% \strut so rows keep an even pitch whether or not the problem has a fraction
% or an exponent in it -- a column of drill lines that breathes unevenly is
% harder to run down with a pencil.
\newcommand{\z}[2]{%
  \stepcounter{zad}%
  \btstoreans{#2}%
  \par\noindent
  \makebox[2.1em][r]{\color{btgrey}\footnotesize\thezad.\ }%
  #1\,\hfill\rule[-0.2ex]{2.6em}{0.4pt}%
  \strut\par\vspace{\btzadskip}%
}

% Wider exercises -- a word puzzle, a sequence with many terms -- that will not
% sit three to a row. Same numbering and the same answer store; only the
% surrounding column count differs, and that is the set's business.
\newcommand{\zz}[2]{%
  \stepcounter{zad}%
  \btstoreans{#2}%
  \par\noindent
  \makebox[2.1em][r]{\color{btgrey}\footnotesize\thezad.\ }%
  % The parbox must leave room for EVERYTHING beside it: the 2.1em number box
  % before it and the 2.6em answer rule after it, plus a gap. Reserving 3.6em
  % for 4.7em of furniture overflowed every wide line by exactly the missing
  % 1.1em -- 12.045pt at 11pt, and the constant width across sixteen boxes is
  % what said it was arithmetic rather than bad line breaking.
  \parbox[t]{\dimexpr\linewidth-5.6em\relax}{\raggedright #1\strut}%
  \hfill\rule[-0.2ex]{2.6em}{0.4pt}%
  \par\vspace{\btzadskip}%
}

% ------------------------------------------------------------
%  The set
% ------------------------------------------------------------
% \begin{zestaw}{title}{stars}{target time}{columns}
%
% The head carries the three things that are true of every exercise below it --
% what kind, how hard, how long the whole set should take -- and the foot
% carries the two the reader writes in. Those two boxes are the point of the
% book: a set you did not time and did not score is just a page of sums.
\newcommand{\DifficultyStars}[1]{%
  \ifcase#1
  \or {\color{orange}\faStar}\faStar[regular]\faStar[regular]%
  \or {\color{orange}\faStar\faStar}\faStar[regular]%
  \or {\color{orange}\faStar\faStar\faStar}%
  \fi
}

\newcommand{\btsetfoot}{%
  \par\vspace{6pt}%
  \noindent\textcolor{btrule}{\rule{\linewidth}{0.5pt}}\par\vspace{4pt}%
  \noindent
  {\bfseries Czas:}~\rule[-0.4ex]{6em}{0.4pt}\hfill
  {\bfseries Poprawne:}~\rule[-0.4ex]{3em}{0.4pt}\,/\,\thezad\hfill
  {\bfseries Data:}~\rule[-0.4ex]{6em}{0.4pt}%
  \par
}

\newenvironment{zestaw}[4]{%
  \refstepcounter{zestaw}%
  \par\addvspace{10pt}%
  % Head
  \noindent\textcolor{btink}{\rule{\linewidth}{1.2pt}}\par\vspace{3pt}%
  \noindent
  {\large\bfseries ZESTAW \thezestaw}\quad{\large #1}%
  \hfill \DifficultyStars{#2}\quad
  \tikz[baseline=(b.base)]{\node[fill=black!8, rounded corners=3pt,
    inner xsep=5pt, inner ysep=2.5pt, outer sep=0pt,
    font=\small\bfseries] (b) {\faClock[regular]\ cel: #3};}%
  \par\vspace{2pt}%
  \noindent\textcolor{btrule}{\rule{\linewidth}{0.5pt}}\par\vspace{7pt}%
  \addcontentsline{toc}{section}{Zestaw \thezestaw\ --- #1}%
  % multicol warns at one column, and it is right to: a one-column set is just
  % a list. The wide sets (word puzzles, long prompts) are exactly that.
  \def\btcols{#4}%
  \ifnum#4>1 \begin{multicols}{#4}\raggedcolumns \fi
}{%
  \ifnum\btcols>1 \end{multicols}\fi
  \btsetfoot
}

% ------------------------------------------------------------
%  The answers appendix
% ------------------------------------------------------------
% Every set's answers, in order, at the back. \ifcsdef rather than an assumed
% range: a set that stored nothing prints nothing, instead of an empty heading
% that reads like a missing page.
\newcounter{btansi}
\newcommand{\odpowiedzi}{%
  \chapter*{Odpowiedzi}%
  \addcontentsline{toc}{chapter}{Odpowiedzi}%
  \markboth{Odpowiedzi}{Odpowiedzi}%
  \small\raggedright
  \setlength{\parindent}{0pt}%
  \setcounter{btansi}{0}%
  \loop\ifnum\value{btansi}<\value{zestaw}%
    \stepcounter{btansi}%
    \ifcsdef{btans@\thebtansi}{%
      \par\addvspace{7pt}%
      {\bfseries\color{btink}Zestaw \thebtansi}\par\vspace{1pt}%
      \csuse{btans@\thebtansi}\par
    }{}%
  \repeat
}

% ------------------------------------------------------------
%  Chapter openers and running heads
% ------------------------------------------------------------
\makeatletter
\renewcommand{\@makechapterhead}[1]{%
  \vspace*{4pt}%
  {\parindent\z@\raggedright\normalfont
   {\small\bfseries\color{btgrey}\MakeUppercase{\chaptername\ \thechapter}}%
   \par\vspace{2pt}%
   {\huge\bfseries #1}%
   \par\vspace{5pt}%
   \textcolor{btink}{\hrule height 1.4pt}%
   \vspace{12pt}%
  }%
}
\renewcommand{\@makeschapterhead}[1]{%
  \vspace*{4pt}%
  {\parindent\z@\raggedright\normalfont
   {\huge\bfseries #1}%
   \par\vspace{5pt}%
   \textcolor{btink}{\hrule height 1.4pt}%
   \vspace{12pt}%
  }%
}
% A blank verso carries no running head. The page \cleardoublepage inserts is
% blank in the BODY only; left alone it ships a head and a page number over an
% otherwise empty sheet.
\renewcommand{\cleardoublepage}{%
  \clearpage
  \if@twoside
    \ifodd\c@page\else
      \hbox{}\thispagestyle{empty}\newpage
      \if@twocolumn\hbox{}\newpage\fi
    \fi
  \fi
}
\makeatother

\pagestyle{fancy}
\fancyhf{}
\fancyhead[LE,RO]{\small\thepage}
\fancyhead[RE,LO]{\small\itshape\nouppercase{\leftmark}}
\renewcommand{\headrulewidth}{0.4pt}

% Must come AFTER \pagestyle{fancy}: fancyhdr installs its own \chaptermark at
% that point and silently discards anything defined earlier. The symptom is a
% running-head change that appears to do nothing at all.
\makeatletter
\renewcommand{\chaptermark}[1]{\markboth{#1}{#1}}
\makeatother
